\section{Sanctity, Legacy, and Unresolved Questions}

Saint Ambrose of Milan did not enter the episcopate as a secluded scholar or a desert ascetic; he arrived at the altar straight from the governor's tribunal, carrying into the Church the habits of Roman jurisprudence, public command, and civic responsibility. What turned those worldly capacities into Christian sanctity was not the preservation of senatorial status, but its deliberate surrender: the rapid divestment of personal wealth, the discipline of learning Scripture on his knees while already preaching daily to his flock, the melting of golden chalices to ransom barbarian captives, the refusal to let armed crowds shed blood for his basilica, and the deep pastoral tenderness that wept beside a brother's bier and reassured a troubled African mother about Saturday fasts.

Across Catholic history, four enduring legacies flow directly from his twenty-three years in Milan. In Trinitarian and sacramental theology, \work{De Fide}, \work{De Spiritu Sancto}, and \work{De Mysteriis} gave the Latin Church clear, scriptural, and pastorally concrete language for the co-equal divinity of Son and Spirit and the real transformation of the Eucharistic elements by the words of Christ. In Christian ethics and ministry, \work{De Officiis} and \work{De Paenitentia} anchored clerical duty in humility, mercy, and service to the poor rather than Stoic pride or Novatianist rigorism. In public worship, his iambic dimeter hymns taught Western laypeople to sing Nicene dogma aloud in the darkness of vigil, inaugurating fifteen centuries of Latin liturgical poetry. And in the relationship between spiritual and civil authority, his refusal to surrender the basilicas to Justina and his private call to Emperor Theodosius after Thessalonica established a lasting moral witness: ``the Emperor is within the Church, not above it'' (\work{Sermon against Auxentius} 36).

Honesty leaves his portrait unvarnished. The same fearless resolve that protected innocent blood after Thessalonica and pleaded clemency for Eugenius's beaten followers also drove Ambrose at Callinicum in 388 to shield Christian rioters from imperial justice and to deploy harsh anti-Jewish rhetoric against a destroyed synagogue. Reverence for a Father and Doctor of the Church does not require defending every political intervention or polemical outburst of fourth-century Roman Christendom; rather, it measures his whole life by the Gospel he preached and the mercy he invoked on his deathbed: ``we have a good Lord.''

Several historical questions remain open in the surviving evidence. Ambrose's exact birth year (c.~339--340 versus c.~333) and birthplace (probable Trier) cannot be demonstrated from his own writings. Whether his episcopal consecration took place on 7~December 374 or 7~December 373 remains debated in modern chronology. The exact date of Satyrus's death (378 versus 379, September versus October) diverges between historical reconstructions and liturgical calendars. The full extent of Ambrose's authentic hymnal beyond the four externally corroborated hymns, the stenographic genesis of \work{De Sacramentis}, and the precise relationship between Paulinus's fifth-century hagiographic vignettes and unrecorded fourth-century events continue to require critical discrimination. Yet even after every late legend is set in its proper register, the historical figure who remains---magistrate turned pastor, grieving brother, patient teacher of Augustine, poet of the dawn, and fearless guardian of the altar---stands in full stature among the foundational Fathers of the West.
